\documentclass[12pt,a4paper]{article}
\usepackage[margin=2.5cm]{geometry}
\usepackage{times}
\usepackage[utf8]{inputenc}
\usepackage[T1]{fontenc}
\usepackage{amsmath,amssymb}
\usepackage{graphicx}
\usepackage{booktabs}
\usepackage{hyperref}
\usepackage{cite}
\usepackage{float}
\usepackage{setspace}

\onehalfspacing

\title{\textbf{Your EI Journal Paper Title}}

\author{
  First Author\textsuperscript{1},
  Second Author\textsuperscript{2},
  Third Author\textsuperscript{1,*} \\[10pt]
  \textsuperscript{1}School of Engineering, University One, City, Country \\
  \textsuperscript{2}School of Engineering, University Two, City, Country \\[6pt]
  \textsuperscript{*}Corresponding author. Email: third@example.com
}

\date{}

\begin{document}
\maketitle

\begin{abstract}
Provide a concise summary of the engineering research including background, methodology, results, and conclusions. Typically 150-300 words.
\end{abstract}

\noindent\textbf{Keywords:} keyword one, keyword two, keyword three, keyword four

\section{Introduction}
Introduce the engineering problem and research objectives.

\section{Theoretical Background}
Present the theoretical framework.

\section{Methodology}
\subsection{Experimental Setup}
Describe the experimental or simulation setup.

\subsection{Measurement Procedures}
Detail measurement methods and parameters.

\section{Results and Discussion}
\subsection{Experimental Results}
Present results with figures and tables.

\subsection{Analysis}
Analyze and discuss the results.

\subsection{Comparison with Existing Methods}
Compare with prior approaches.

\section{Conclusion}
Summarize the main contributions and engineering implications.

\section*{Acknowledgments}
This work was supported by ...

\bibliographystyle{unsrt}
\bibliography{references}

\end{document}
